\chapter{Aphasia (Difficulty Speaking)}

\section*{Definition}
Aphasia is an acquired disorder of language that can affect speaking, understanding, reading, or writing. It usually follows injury to language networks in the dominant hemisphere, often the left, but language is distributed across several connected regions. Aphasia differs from dysarthria, which is impaired motor control of speech, and from apraxia of speech, which is difficulty planning speech movements.

\section*{When to See a Doctor}

\begin{itemize}
 \item Sudden onset of aphasia or difficulty understanding speech -- treat it as a stroke emergency, even if it improves
 \item Aphasia with facial droop, arm weakness, or confusion (stroke symptoms)
 \item Progressive language decline over months (degenerative dementia)
 \item Transient aphasia with headache or visual changes; migraine is one possibility, but stroke and seizure must be excluded
 \item Aphasia after head trauma
 \item Any new language difficulty requires prompt medical evaluation
\end{itemize}

\section*{How It Develops}
Language depends on connected networks in the dominant frontal, temporal, parietal, and subcortical regions. Injury can selectively affect word retrieval, grammar, understanding, repetition, reading, or writing. Traditional labels such as Broca, Wernicke, global, and conduction aphasia can describe patterns, but real-world aphasia is often mixed and does not map neatly to one brain location.

\section*{Causes}
\begin{itemize}
 \item Stroke (ischemic or hemorrhagic -- most common cause)
 \item Transient ischemic attack or minor stroke causing temporary language symptoms
 \item Traumatic brain injury
 \item Brain tumor (primary or metastatic)
 \item Brain infection (encephalitis, abscess)
 \item Dementia: Alzheimer disease, primary progressive aphasia, frontotemporal dementia
 \item Seizure (Todd paresis -- transient post-ictal aphasia)
 \item Migraine with aura (transient aphasic symptoms)
 \item Metabolic encephalopathy
 \item Medications: Sedatives, anticonvulsants (reversible)
\end{itemize}

\section*{Related Symptoms}
\begin{itemize}
 \item Dysarthria
 \item Dysphagia
 \item Hemiparesis (one-sided weakness)
 \item Hemisensory loss
 \item Visual field deficits
 \item Cognitive impairment
 \item Apraxia
 \item Memory loss
 \item Depression
\end{itemize}

\section*{Medical Evaluation}
\begin{itemize}
 \item Neurological examination including full cranial nerve assessment.
 \item Urgent brain imaging, usually noncontrast CT initially when rapid stroke or bleeding assessment is needed; MRI may provide more detail when available and appropriate.
 \item Formal language assessment by speech-language pathologist (Boston Diagnostic Aphasia Examination, Western Aphasia Battery).
 \item Assessment of other cognitive domains (memory, attention, executive function).
 \item EEG if seizure suspected.
 \item Vascular imaging of the head and neck when indicated by the suspected stroke mechanism; carotid imaging is not required for every aphasia presentation.
\end{itemize}

\section*{Treatment Options}
\begin{itemize}
 \item Acute treatment follows the cause and time window: eligible ischemic strokes may receive reperfusion treatment, while hemorrhage, seizure, infection, or tumor requires cause-specific care in a specialist setting.
 \item Speech-language therapy should be individualized and may include restorative language practice, supported communication, and compensatory strategies. Improvement can occur in both early and chronic stages.
 \item Pharmacotherapy for specific etiologies.
 \item Group therapy and communication partner training.
 \item Augmentative and alternative communication (AAC) for severe persistent aphasia.
 \item Treatment of depression (common after aphasia).
\end{itemize}

\section*{Self-Care and Home Management}
\begin{itemize}
 \item For acute onset, call emergency services immediately and record the time the person was last known well. Do not drive the person if an ambulance is available.
 \item Use simple, short sentences and yes/no questions.
 \item Allow extra time for response; do not finish the person's sentences.
 \item Use gestures, pictures, or writing to support communication.
 \item Maintain a calm environment with minimal distractions.
 \item Seek speech-language pathology evaluation.
 \item Do not self-treat; underlying cause requires urgent diagnosis.
\end{itemize}

\section*{References}
\begin{thebibliography}{99}
\bibitem{aphasia:ref1} Berthier ML. ``Poststroke aphasia: epidemiology, pathophysiology and treatment.'' \textit{Drugs Aging}. 2005;22(2):163-182.
\bibitem{aphasia:ref2} Brady MC, Kelly H, et al. ``Speech and language therapy for aphasia following stroke.'' \textit{Cochrane Database Syst Rev}. 2016;(6):CD000425.
\bibitem{aphasia:ref3} Kirshner HS. ``Aphasia and aphasia syndromes.'' In: \textit{Bradley\'s Neurology in Clinical Practice}, 8th ed. Elsevier, 2021.
\bibitem{aphasia:ref4} National Institute on Deafness and Other Communication Disorders. ``Aphasia.'' NIH Pub. No. 97-4257, 2015.
\bibitem{aphasia:ref5} UpToDate. ``Approach to the patient with aphasia.'' Wolters Kluwer, 2023.
\bibitem{aphasia:ref6} National Institute on Deafness and Other Communication Disorders. Aphasia. Updated 2025.
\bibitem{aphasia:ref7} American Heart Association/American Stroke Association. Guidelines for adult stroke rehabilitation and recovery. 2016; updated stroke guidance 2026.
\end{thebibliography}
